
\chapter{Ejemplos y Aplicaciones}

\begin{verbatim}
    
 _____
|___  |
   / /
  / /
 /_/

 _____  _                      _
| ____|(_) ___ _ __ ___  _ __ | | ___  ___
|  _|  | |/ _ \ '_ ` _ \| '_ \| |/ _ \/ __|
| |___ | |  __/ | | | | | |_) | | (_) \__ \
|_____|/ |\___|_| |_| |_| .__/|_|\___/|___/
     |__/               |_|
            _          _ _                _
 _   _     / \   _ __ | (_) ___ __ _  ___(_) ___  _ __   ___  ___
| | | |   / _ \ | '_ \| | |/ __/ _` |/ __| |/ _ \| '_ \ / _ \/ __|
| |_| |  / ___ \| |_) | | | (_| (_| | (__| | (_) | | | |  __/\__ \
 \__, | /_/   \_\ .__/|_|_|\___\__,_|\___|_|\___/|_| |_|\___||___/
 |___/          |_|




Contenido
=========

1. Introducción.
2. Probabilidades y Monedas.
3. Juguemos con Dados.
4. El Valor de Pi.
5. Transportes y Mármol.
6. Un Sistema de Inventarios.
7. Sistemas de Colas.
8. Resumen.
9. Ejercicios.




1. Introducción.
================

En los capítulos anteriores se han cubierto todas las herramientas teóricas que se necesitan para construir un modelo de simulación. Podemos presentar ahora ejemplos, aplicaciones y análisis de estos sistemas.

Para cada ejemplo que se presente se tomará en cuenta la dificultad de resolver dicho modelo mediante un procedimiento analítico matemático, si es que este existe. Asímismo se debe considerar las posibilidades prácticas y de costo de manipular los problemas en la vida real, contra la manipulación del modelo digital.

En la vida real los modelos de simulación suelen utilizarse para diferentes fines, tales como el diseño de nuevos sistemas, el análisis de la eficiencia de sistemas reales, el análisis para realizar cambios a sistemas existentes y para entrenar personas en el uso y peligros de sistemas. Finalmente una nueva aplicación ha surgido en los últimos años que consiste en la aplicación de simulación a la teoría de juegos y al uso de juegos por computadora.

Se presentan algunos ejemplos que sirven como prototipos. Cualquier nuevo modelo de simulación puede partir de uno de los prototipos aquí presentados y luego realizar los cambios, adaptaciones y extensiones necesarias.




2. Probabilidades y Monedas.
============================

Se apuesta con una moneda legal y se tiene una cantidad de dinero inicial de $30. La apuesta inicial será de $10. En general la ganacia es igual a la cantidad apostada, en caso de perder se pierde la cantidad apostada.

El juego consiste en tirar la moneda, si cae escudo se gana y si cae cara se pierde. Para jugar se llevará a cabo la siguiente estrategia, la apuesta inicial es de $10. Si se gana se vuelven a apostar $10, si se pierde se apuesta el doble de la vez anterior. De esta forma si apuesto $10 y pierdo la siguiente vez se apostarán $20, si pierdo la siguiente vez se apostarán $40. Si no se tiene suficiente dinero para cubrir la apuesta se apostará todo lo que se tiene. El juego continuará hasta haber ganado $50 o bien haber perdido todo y quedar en quiebra.

Otra forma de describir el juego es decir que se inicia con una apuesta de x=$10, si pierdo apuesto 2*x, si vuelvo a perder 4*x, y así sucesivamente. Si gano apuesto "x". Al igual que antes el juego continua hasta obtener $50 o quedar con $0 en estado de quiebra.

¿Cuál es la probabilidad de ganar $50 con este juego?



Solución.
---------


Evidentemente aunque el problema es de un planteamiento sencillo su solución analítica es difícil. Por eso debemos requirir a un mecanismo de simulación.

Para simular la moneda se utilizará la transformada inversa una variable aleatoria uniforme. La variable "r" es un valor aletorio en [0,1], "e" significa el evento que la moneda caiga escudo y "c" significa que caiga una corona.

---------------
x   p(x)  P(x)
---------------
e   0.50  0.50
c   0.50  1.00
---------------

 
De esta forma el algoritmo de generación de números aleatorios es:

Se escoge un valor de r en el intervalo [0,1[
Si r está en [0.00 , 0.50[ se gana, es equivalente a un escudo.
Si r está en [0.50 , 1.00[ se pierde, es equivalente a cara.



A continuación se presentan cinco corridas manuales de este problema para estudiar su comportamiento.


--------------------------------------------------
     Dinero                     Dinero
No.  Inicial  Apuesta   r       Final     Meta
--------------------------------------------------
 1     30       10     0.0399     40      no
       40       10     0.4885     50      si**

 2     30       10     0.1825     40      no
       40       10     0.3264     50      si**

 3     30       10     0.7124     20      no
       20       20     0.2524     40      no
       40       10     0.6140     30      no
       30       20     0.8872     10      no
       10       10     0.6564      0      quiebra

 4     30       10     0.0370     40      no
       40       10     0.8086     30      no
       30       20     0.1650     50      si**

 5     30       10     0.5582     20      no
       20       20     0.9195      0      quiebra
--------------------------------------------------



Si estas cinco corridas fueran significativas del problema bajo estudio se podría decir que el juego se gana con una probabilidad de 0.60 de las ocasiones y se pierde con una probabilidad de 0.40. Esto ya que en la simulación se ganó el juego en 3 de las 5 corridas y se perdió en 2 de las 5 corridas. Sin embargo 5 corridas NO son suficientes para establecer este número.

Todo proceso de simulación pasa por una fase conocida como fase transciente. Durante este período los resultados obtenidos no son significativos del experimento bajo estudio. Cuando se han realizado un número significativamente grande de corridas se puede decir que el problema se encuentra en una fase estable. Durante la fase estable se pueden obtener los resultados deseados.

A continuación se presenta una tabla donde se indica el número de corridas, o números de veces que se ejecuta el juego, y de este total de corridas el número de veces que se logra la meta de obtener $50.

-----------------------
        No.    Corridas
    Corridas   Exitosas
-----------------------
         10           6
        100          67
      1,000         617
     10,000       6,106
    100,000      61,059
  1,000,000     609,825
 10,000,000   6,093,850
100,000,000  60,943,397
-----------------------


Los resultados de los procesos de simulación no se suelen presentar de forma puntual, sino, que dentro de un intervalo estadístico de confianza. 

A continuación se presenta un ejemplo con 10000 corridas del problema. Para lograr el cálculo adecuado del intervalo de confianza se supondrá que cada vez que se gana el juego la simulación produce un 1 y cada vez que se llega a un estado de quiebra la simulación produce un valor de 0. Asímismo se utilizará un intervalo de confianza de 0.95.


Para calcular el valor del promedio de éxitos se utilizará la fórmula:

_      /  6106  \
x  =  (  ------  )  =  0.6106
       \ 10000  /



Ahora bien, para calcular el valor de la variancia y de la desviación estándar procedemos a utilizar la siguiente fórmula. Se supondrá que el valor de 1 se utiliza cuando se gana y el valor de 0 se utiliza cuando se se queda en bancarrota.

         __
         \       _ 
         /_ (x - x)^2
s^2  =   ------------
            (n - 1)


                         
         6106 * (1 - 0.6106)^2   +  3894 * (0 - 0.6106)^2
s^2  =   ------------------------------------------------
                               9999



        925.8672 + 1451.8092
s^2  =  --------------------
               9999



s^2  =  0.2378



Una vez que se ha obtenido el valor de la variancia, la desviación estándar se obtiene al sacar la raíz cuadrada de la misma.


s^2  =   0.2378


s  =   sqrt(0.2378)


s  =   0.4876


Procedemos ahora a calcular la desviación estándar del promedio y así hacer el cálculo del intervalo de confianza.

  _     
s(x)  =  s / sqrt(n)

  _
s(x)  =  0.4876 / sqrt(10000)

  _
s(x)  =  0.4876 / 100

  _
s(x)  = 0.0049



Para calcular el intervalo de confianza se construyen los valores del límite inferior y el límite superior. Para ello se utiliza un valor de z0 que abarque el 0.95 de la distribución normal.


|
|
|                         . ' .
|                       .       .
|                     .           .
|            |      .               .      |
|            |    .                   .    |
|            | .                         . |
|          . |                             | .
|     .      |                             |      .
|____________|_____________________________|________________
           -1.96                          1.96



Finalmente para calcular los intervalos de confianza utilizamos las fórmulas de siempre:


       _                  _
Li  =  x     -  z0   *  s(x)

Li  = 0.6106 - 1.96 * 0.0049

Li  = 0.6010


       _                 _
Ls  =  x     +  z0  *  s(x)
                        
Ls  = 0.6106 + 1.96 * 0.0049

Ls  = 0.6202



De esta forma se puede concluir con una confianza del 95% que el verdadero valor de esta simulación se encuentra en el intervalos compuesto por [0.6010 , 0.6202].




3. Juguemos con Dados.
======================

El juego de casino de los dados conocido como "craps" requiere que un jugador lance dos dados. Para el juego principal, pues existen muchas otras derivaciones, se siguen las siguientes reglas.

() Si en el primer tiro el jugador consigue un 7 o un 11 gana.

() Si en el primer tiro el jugador consigue un 2,3 o un 12 pierde.

() Si en el primer tiro se consigue un 4,5,6,8,9,10 el jugador debe continuar tirando hasta que ocurra uno de dos eventos. Si consigue repetir el valor inicial gana, si saca un 7 en cualquier otro momento pierde.

Se desea determinar la probabilidad de ganar en este juego.



Solución.
---------



A continuación se presentan algunas corridas de ejemplo del juego.

----------------------------
No.    Jugadas     Resultado
----------------------------
1.   7               gana

2    11              gana

3    4,6,11,8,4      gana

4    6,11,8,7        pierde

4    12              pierde

5    10,8,6,7        pierde

6    8,6,9,5,7,8     gana
----------------------------


Para poder simular este juego adecuadamente, un dado que se tira dos veces, dos dados, o bien la suma de dos dados perfectos.


Para simular un dado se requiere una distribución uniforme discreta entre [1,6]. De esta forma se generan dos números aleatorios y se suman. A continuación se presenta el proceso mediante la transformada inversa.

---------------
x   p(x)  P(x)
---------------
1   1/6   1/6
2   1/6   2/6
3   1/6   3/6
4   1/6   4/6
5   1/6   5/6
6   1/6   6/6
--------------


Con base en la tabla anterior se genera el siguiente algoritmo:

Se genera un valor de "r" entre [0,1[.

Si r está entre [0   , 1/6[ entonces x = 1
Si r está entre [1/6 , 2/6[ entonces x = 2
Si r está entre [2/6 , 3/6[ entonces x = 3
Si r está entre [3/6 , 4/6[ entonces x = 4
Si r está entre [4/6 , 5/6[ entonces x = 5
Si r está entre [5/6 , 6/6[ entonces x = 6


La otra forma es simular la distribución de la suma, la cual se puede comprobar mediante la siguiente tabla.

+-----+-------------------------+
|     |  1   2   3   4   5   6  |
+-----+-------------------------+
|  1  |  2   3   4   5   6   7  |
|     |                         |
|  2  |  3   4   5   6   7   8  |
|     |                         |
|  3  |  4   5   6   7   8   9  |
|     |                         |
|  4  |  5   6   7   8   9  10  |
|     |                         |
|  5  |  6   7   8   9  10  11  |
|     |                         |
|  6  |  7   8   9  10  11  12  |
+-----+-------------------------+


Con la tabla anterior se pueden obtener las probabilidades de cada uno de los resultados, los cuales se muestran a continuación.

-------------------
x    p(x)     P(x)
-------------------
2    1/36    1/36
3    2/36    3/36
4    3/46    6/36
5    4/36   10/36
6    5/36   15/36
7    6/36   21/36
8    5/36   26/36
9    4/36   30/36
10   3/36   33/36
11   2/36   35/36
12   1/36   36/36
-------------------

Con base en la tabla anterior se construye el siguiente algoritmo:

Se escoge un valor de r en [0,1[
Si r esta entre [ 0/36 ,  1/36[ entonces x = 2 
Si r esta entre [ 1/36 ,  3/36[ entonces x = 3 
Si r esta entre [ 3/36 ,  6/36[ entonces x = 4 
Si r esta entre [ 6/36 , 10/36[ entonces x = 5 
Si r esta entre [10/36 , 15/36[ entonces x = 6 
Si r esta entre [15/36 , 21/36[ entonces x = 7 
Si r esta entre [21/36 , 26/36[ entonces x = 8 
Si r esta entre [26/36 , 30/36[ entonces x = 9 
Si r esta entre [30/36 , 33/36[ entonces x = 10
Si r esta entre [33/36 , 35/36[ entonces x = 11
Si r esta entre [35/36 , 36/36[ entonces x = 12


La simulación se puede realizar de la siguiente manera:

Se tiran dos dados y se obtiene su suma.
Si es 7,11 gana.
Si es 2,3,12 pierde.
Para cualquier otro valor se establece como meta, se sigue tirando el dado hasta que ocurra uno de los dos eventos siguientes: si sale un 7 se pierde, si sale la meta se gana.



A continuación se presenta una tabla con simulaciones simultáneas de este juego. Entre mayor sea el número de veces que se ha ejecutado el programa, más se acerca el resultado obtenido al resultado real del juego.

-----------------------
        No.    Corridas
    Corridas   Exitosas
-----------------------
         10           3
        100          51
      1,000         503
     10,000       4,885
    100,000      49,818
  1,000,000     491,823
 10,000,000   4,928,861
100,000,000  49,309,658
-----------------------


Si el proceso se simula para un número lo suficientemente grande de jugadas, la proporción del número de juegos ganados tendrá una distribución normal con las siguientes características:

_
x  =  0.493

  _
s(x)  =  0.5 / sqrt(n)



    

4. El Valor de Pi.
==================

A continuación se calculará una aproximación al valor de pi utilizando un proceso geométrico. Dada la siguiente figura, se pueden estimar las áreas del cuadrado de base 1 y del círculo de radio 1. Se desea usar esta información para encontrar una aproximación al valor de pi.


r2
 |
 |
1|________________
 |  .             |
 |       .        |
 |          .     |
 |            .   |
 |                |
 |              . |
 |                |
 |               .|
 |________________|_____r1
                  1



Solución.
---------


El área del cuadrado es fácil de calcular pues consiste en múltiplicar dos de sus lados o aristas.

área_cuadrado = lado * lado

área_cuadrado = 1 * 1

área_cuadrado = 1



El área del círculo se puede estimar mediante la fórmula:

                __      
área_círculo =  ||  *  r^2


Como solamente contamos con la cuarta parte de un círculo debemos dividir el valor entre 4:

                __      
área_círculo =  ||  *  r^2   /  4

                __      
área_círculo =  ||  *  1^2   /  4

                __     
área_círculo =  ||  /  4



De tal forma que que si se escoge un punto de forma aleatoria en el cuadrado de área 1, la probabilidad que este punto caiga dentro del área del círculo va a estar dada por:


                                  área_círculo
p([r1,r2] dentro del círculo)  =  -------------
                                  área_cuadrado


                                     __
                                     ||  /  4
p([r1,r2] dentro del círculo)  =  -------------
                                        1


                                     __
p([r1,r2] dentro del círculo)  =     ||  /  4



Si calculamos un número de puntos lo suficientemente grande, el valor del círculo va a estar dado por la proporción  x/n  donde "x" representa el número de puntos dentro del círculo y "n" el total de puntos generados.



Es importante observar que :

                  __
           x      ||
 Límite   ---  =  --
n -> +oo   n       4



En consecuencia para un número grande de corridas se puede suponer que:

        __
 x      ||
---  =  --
 n       4


__     4*x
||   = ---
        n


Para determinar si un punto se encuentra dentro del círculo utilizamos la ecuación que define la circunferencia del círculo para luego encontrar su interior.

x^2 + y^2  =  r^2

sqrt(x^2 + y^2) = r



Si se usan dos números r1 y r2 generados aleatoriamente en [0,1[ se tiene que:

sqrt(r1^2 + r2^2) = r

sqrt(r1^2 + r2^2) = 1

sqrt(r1^2 + r2^2) <= 1




Utilizando la ecuación anterior, el algoritmo para descubrir el valor de pi sería el siguiente.

() Generar dos números uniformes en [0,1[ que denominaremos r1 y r2.

() Evaluar si estos dos números se encuentran dentro del círculo mediante la fórmula: 
sqrt(r1^2 + r2^2) <= 1

() Repetir los pasos anteriores hasta obtener "n" corridas, "x" representa el número de puntos dentro del círculo. En ese momento el valor de pi estará dado por:
__
||  = 4 * (x/n)




A continuación se muestra una tabla con unas cuantas ejecuciones de este algoritmo.

---------------------------------------------------
                       sqrt( r1^2     ¿Dentro del
No.     r1       r2        + r2^2 )    círculo? 
---------------------------------------------------
 1    0.4597   0.6449     0.7920       si
 2    0.8029   0.2382     0.8375       si
 3    0.1347   0.3156     0.3431       si
 4    0.8140   0.8111     1.1491       no
 5    0.0358   0.8024     0.8032       si
---------------------------------------------------


Por lo tanto el valor estimado de pi estará dado por:

__
||  = 4 * (x/n)

__
||  = 4 * (4/5)

__
||  = 16/5

__
||  = 3.20



Como es evidente con solamente 5 corridas el valor estimado de pi difiere en mucho de su valor real. Por lo tanto si se desea aproximar más el valor de la estimación a su valor verdadero es necesario aumentar el número de corridas de forma significativa. Cabe entonces la pregunta de cuántas corridas es necesario realizar para que la estimación de pi difiera de su valor verdadero en una cantidad de 0.001 con un nivel de confianza del 95%. 

Esta pregunta se puede plantear como un problema de intervalos de confianza. El valor de pi asterisco es la aproximación a pi.


  / __              __*      __         \
p(  || - 0.001  <=  ||   <=  || + 0.001  )  <=  0.95
  \                                     /



  / __              4*x      __         \
p(  || - 0.001  <=  ---  <=  || + 0.001  )  <=  0.95
  \                  n                  /


     __                       __         
  /  || - 0.001       x       || + 0.001 \
p(   ----------  <=  ---  <=  ----------  )  <=  0.95
  \     4             n          4       /



Ahora bien "x" representa la cantidad de puntos que cayeron adentro del cuarto de círculo. Se puede pensar que caer adentro es un éxito y caer afuera es un fracaso, de esta forma "x" sigue una distribución binomial. Como sabemos la distribución binomial tiene la siguiente media y variancia.


E(x)  =  n * p

               __
E(x)  =  n * ( || / 4 )
                
               __
E(x)  =  n * ( || / 4 )



Y se puede deducir que:


  / x \     1
E(  -  ) =  -  *  E(x)
  \ n /     n


                        __
  / x \     1           ||
E(  -  ) =  -  *  n  * ----
  \ n /     n            4


             __
  / x \      ||
E(  -  ) =  ----
  \ n /       4



   / x \       __
 E(  -  ) =  ( || / 4 )
   \ n /   




V(x)  =  n * p * q


                  __               __
               /  ||  \     /      ||  \ 
V(x)  =  n *  (  ----  ) * (  1 - ----  )
               \   4  /     \       4  / 


               __               __
V(x)  =  n * ( || / 4 ) * ( 1 - || / 4 )





  /  x  \      /  1  \ 
V(  ---  )  = (  ---  ) *  V(x)
  \  n  /      \ n^2 / 



                                    __               __   
  /  x  \       /  1  \          /  ||  \     /      ||  \ 
V(  ---  )  =  (  ---  ) * n *  (  ----  ) * (  1 - ----  )
  \  n  /       \ n^2 /          \   4  /     \       4  / 



                                __                __   
  /  x  \       /  1  \      /  ||  \      /      ||  \ 
V(  ---  )  =  (  ---  ) *  (  ----  )  * (  1 - ----  )
  \  n  /       \  n  /      \   4  /      \       4  / 



  /  x  \                  __               __
V(  ---  )  =  (1 / n) * ( || / 4 ) * ( 1 - || / 4 )
  \  n  /    



                   ________________________________
   /  x  \        /          __             __
de(  ---  )  =  \/ (1/n) * ( ||/4 ) * ( 1 - ||/4 )
   \  n  / 




Como el valor de "n" es lo suficientemente grande, la distribución binomial se puede calcular por medio de la distribución normal, utlizando la misma esperanza y la misma variacia. En general, al normalizar se utiliza la fórmula:


      x  -  u
z0 =  --------
         o


Que en nuestro caso se convertiría en:

                  
     (x/n)  -   u(x/n)
z0 = ---------------
          o(x/n)
             



     (x/n) - E(x/n)
z0 = -------------
        de(x/n)




El objetivo es encontrar el número de corridas, lo que corresponde a despejar "n". Por ello resulta más sencillo normalizar el lado derecho.



        __
      ( || + 0.001 )/4  -  E(x/n)
z0 =  ---------------------------
                    de(x/n)



              __              __
              || + 0.001      ||
              ----------   -  --
                 4             4
z0 =  ------------------------------------
          ________________________________
         /          __             __
       \/ (1/n) * ( ||/4 ) * ( 1 - ||/4 )




        ________________________________      __             __
       /          __             __           || + 0.001     ||
z0 * \/ (1/n) * ( ||/4 ) * ( 1 - ||/4 )    =  ----------  -  --
                                                 4            4




        ________________________________      __               __
       /          __             __           ||  +  0.001  -  ||
z0 * \/ (1/n) * ( ||/4 ) * ( 1 - ||/4 )    =  -------------------
                                                      4




        ________________________________
       /          __             __           0.001
z0 * \/ (1/n) * ( ||/4 ) * ( 1 - ||/4 )    =  -----
                                                4



 
                 __             __           0.001^2
z0^2 * (1/n) * ( ||/4 ) * ( 1 - ||/4 )    =  -------
                                               4^2
   




                __             __
4^2 * z0^2  * ( ||/4 ) * ( 1 - ||/4 )
-------------------------------------   =  n
            0.001^2
            



Como se quería un intervalo de 0.95 en la distribución normal se puede sustituir el valor de z0 = 1.96. De igual forma al sustituir el valor de pi se obtiene lo siguiente:


                      __             __
      4^2 * z0^2  * ( ||/4 ) * ( 1 - ||/4 )
n  =  --------------------------------------
                  0.001^2



   
      4^2 *  1.96^2  * (3.1415... / 4) * ( 1 - 3.1415... /4 )
n  =  -------------------------------------------------------
                              0.001^2



n =  10359897.0849


n =  10359897


n =  10,359,897



Lo que significa que se requieren generar 10359897 números para obtener el valor de pi con una diferencia de 0.001 con un 95% de confianza. Sin embargo esta fórmula es poco utilizable ya que estamos calculando el valor de pi y para ello necesitamos utilizar dicho valor. Por lo tanto es usual sustituir este valor por una cota superior.



                      __             __
      4^2 * z0^2  * ( ||/4 ) * ( 1 - ||/4 )
n  =  -------------------------------------
                  0.001^2




        4^2 * z0^2 * 1 * 1
n  =    ------------------
             0.001^2
               



        4^2 * 1.96^2 * 1 * 1
n  =    --------------------
              0.001^2



n  =  61465600

n  =  61,465,600



Finalmente una diferencia de 0.001 entre el estimado y el valor de pi es muy grande. Si se quieren que el valor de pi tenga 4 dígitos correctos se debe pedir una diferencia menor a 0.00001. Esto convertiría la fórmula anterior en:


        4^2 * 1.96^2 * 1 * 1
n  =    --------------------
              0.00001^2


n  =  614656000000


n  =  614,656,000,000





5. Transportes y Mármol.
========================

La empresa Roma Virtual S.A. tiene asignado un camión especial de transporte para sus estatuas de imitación de mármol. Dicho camión transporta diariamente 5 estatuas. El peso de las estatuas tiene la siguiente distribución de probabilidad.



  f(x)
    |
    |
    |
1/20|                  .
    |                .   .
    |              .       .
    |            .           .
    |          .               .
    |        .                   .
    |      .                       .
    |____.___________________________.____ kilos
        190           210           230



Si el camión de reparto tiene una capacidad de 1000 kilos, entonces estamos interesados en averiguar cuál es la probabilidad que el peso de las estatuas exceda la capacidad del camión. 

Se resolverá el presente problema primero en forma analítica para ver la dificultad del proceso y luego utilizando simulación.




Solución Analítica.
-------------------


La distribución de probabilidad fue dada por medio de un gráfico. Debemos convertir dicho gráfico a una fórmula matemática.

        
            /  x - 190
           /   -------    si 190 <= x <= 210
          /      400
         /     
f(x) =  (
         \ 
          \    230 - x
           \   -------    si 210 <= x <= 230
            \    400




Obtenemos ahora la esperanza, la variancia y la desviación estándar de la distribución.

            _
           | )+oo
           |
E(x)   =   |        x * f(x)  dx
           |
         (_|  -oo



           _
          | )210
          |               / x - 190 \
E(x)  =   |        x  *  (  -------  )   dx
          |               \   400   /
        (_|  190



          _
         | )230
         |                / 230 - x \
     +   |         x  *  (  -------  )    dx
         |                \   400   /
       (_|  210




E(x)  =  210






Y la Variancia:

           _
          | )+oo
          |           _ 
V(x)  =   |      (x - x)^2  * f(x)  dx
          |
        (_|  -oo




           _
          | )210
          |                        / x - 190 \
V(x)  =   |       (x - 210)^2  *  (  -------  )  dx
          |                        \   400   /
        (_|  190




          _
         | )230
         |                         / 230 - x \
     +   |        (x - 210)^2  *  (  -------  )  dx
         |                         \   400   /
       (_|  210




V(x)  =  66.6667



Y finalmente la desviación estándar que se obtiene al sacar la raíz cuadrada de la variancia.

             _____       _________
de(x)  =   \/ V(x)  =  \/ 66.6667   =  8.1649



Para encontrar la probabilidad que la suma total de los pesos individuales de cada estatua exceda el máximo soportado por el camión se debe calcular la siguiente probabilidad, donde x(i) representa el peso de la estatua i.



p(  x(1) + x(2) + x(3) + x(4) + x(5)  >=  1000 )



Que es equivalente a:


                                              
  / x(1) + x(2) + x(3) + x(4) + x(5)      1000 \
p(  --------------------------------  >=  ----  )
  \                5                        5  /


   _
p( x  >=  200)



Por el teorema del límite central conocemos que el promedio tiene las siguientes características.

   _       _
E( x )  =  x

   _
V( x )  =   V(x) /  n

    _       
de( x ) =  de(x) / sqrt(n)



Por lo tanto:


   _
E( x )  =  210


   _
V( x )  =  V(x) / n   

        =  66.6667 / 5  

        =  13.3333


    _       
de( x ) =  de(x) / sqrt(n) 
                          
        =  8.1649 / sqrt(5)

        =  3.6515




Finalmente calculamos la probabilidad



   _                /        200 - 210 \
p( x >= 200 )  =  p(  z  >=  ---------  )
                    \         3.6515   /


               =  p( z >=  -2.74 )


               =  1  -  p( z <= 2.74 )


               =  1  -  0.0031


               =  0.9969



R/ Por lo tanto, la probabilidad de que el peso de las estatuas exceda la capacidad del camión es de 0.9969





Solución por Simulación
-----------------------

Para resolver este ejercicio mediante un programa de simulación se deben realizar las siguientes actividades.


() Simular el peso de 5 estatuas. 

() Sumar estos pesos.

() Comparar la suma con el total de 1000 kilos.

() Repetir el proceso tantas veces como sea necesario.



Para simular el peso de las estatuas se necesita un generador de números aleatorios de acuerdo a la distribución. Para ello se utilizará el método de la transformada inversa.



Para el caso que el valor de x se encuentre en [190,210] calculamos el valor de la siguiente integral.


          _
         | )x
         |       t - 190
F(x)  =  |       -------   dt
         |         400
       (_|  190




Para el caso que el valor de x se encuentre en [210,230] calculamos el valor de la siguiente integral.


                       _
                      | )x
                      |       230 - t
F(x)  =  F(210)  +    |       -------   dt
                      |         400
                    (_|  210




Por lo que la función acumulada tendrá el siguiente comportamiento.


                
              /  0                    si x <= 190
             /   
            /            
           /     (x - 190)^2
          /      ---------             si 190 <= x <= 210
         /          800
        /        
F(x) = (         
        \                       
         \             (x - 230)^2
          \      1  -  -----------    si 210 <= x <= 230
           \              800
            \    
             \   
              \  1                    si x > 230




Ahora para obtener la transformada inversa se toma cada una de las partes de la distribución acumulada. Para la primera parte tenemos.

         
(x - 190)^2
-----------  =  r
     800

         
(x - 190)^2  =  800 * r

                    
(x - 190)  =  +/-  sqrt(800 * r)

                
(x - 190)  =  sqrt(800 * r)

        
 x = 190  +  sqrt(800 * r)




Para la segunda parte de la transformada inversa utilizamos un proceso similar.

               
      (x - 230)^2
1  -  -----------  =  r
         800

         
(x - 230)^2
-----------   =  1 - r
   800



(x - 230)^2  =  800*(1 - r)


                     
(x - 230)  =  +/- sqrt(800 * (1 - r))


x  =  230  -  sqrt(800 * (1 - r))



Para simular el peso de las estatuas se utilizará el siguiente algoritmo:


() Generar un número uniforme "r".


() Si r <= 0.50, entonces el valor de x estará en [190,210]
                   
   x  =  190  +  sqrt(800 * r)


() Si r >= 0.50, entonces el valor de x estará en [210,230]
                 
   x  =  230  -  sqrt(800 * (1 - r))




A continuación se muestra una corrida de este proceso de simulación. 


----------------------------------------------
  No     r         x        Total      Excede?
----------------------------------------------
   1   0.3987   207.8590
   2   0.6398   213.0240
   3   0.3905   207.6749
   4   0.3582   206.9291
   5   0.5285   210.5794   1046.0665   si
----------------------------------------------
   1   0.4744   209.4821
   2   0.6802   214.0043
   3   0.1730   201.7641
   4   0.1774   201.9146
   5   0.9265   222.3320   1049.4972   si
----------------------------------------------
   1   0.2491   204.1159
   2   0.3762   207.3487
   3   0.0855   198.2700
   4   0.8179   217.9292
   5   0.5064   210.1282   1037.7921   si
----------------------------------------------
   1   0.7590   216.1148
   2   0.3782   207.3943
   3   0.5211   210.4274
   4   0.3430   206.5639
   5   0.3604   206.9789   1047.4793   si
----------------------------------------------
   1   0.3672   207.1385
   2   0.0773   197.8657
   3   0.7380   215.5229
   4   0.2600   204.4228
   5   0.8655   219.6288   1044.5786   si
----------------------------------------------




¿Se debe adquirir o no un nuevo camión?
----------------------------------------

Al igual que antes ya sea por medios analíticos o por simulación, se puede analizar la conveniencia de adquirir un nuevo camión.

Se supondrá que se realiza solamente un embarque por día. Cuando la capacidad del camión no es suficiente se envían las estatuas que sobrepasan el peso a una compañía de entregas con un costo adicional de $200 por estatua. 

El costo promedio anual, pago del préstamo por la compra, más manteniniento y seguros de un nuevo camión es de $51900. Si se trabajan 5 días a la semana y 52 semanas al año, ¿cuál alternativa es mejor?



Analicemos la situación anterior de forma analítica.
----------------------------------------------------

El costo de enviar mercadería a través de la otra compañía es:

$200 por estatua

0.9969 de probabilidad por embarque diario.

5 días por semana

52 semanas por año.


Poniendo la anterior junto, se obtiene el costo anual que estará dado por:

($200) * (0.9969) * (5) * (52)  =  $51840


Como 51840 es menor que 51900, parece que no es necesario comprar un nuevo camión si las circunstancias se mantienen así.



Si analizamos la situación mediante simulación, tendríamos que realizar el siguiente proceso.

() Simular un envío de 5 estatuas.

() Si el peso de las 5 estatuas es mayor que 1000, se acumulan $200
   en otro caso no se acumula nada.

() Repetimos el proceso anterior para 5 * 52 = 260 días 
   y determininamos el costo total.


Para que el proceso anterior funcione correctamente hay que ejecutarlo una gran cantidad de veces. Como en cada ocasión producirá un costo diferente se debe sacar un promedio del costo y un intervalo de confianza. Asímismo es natural suponer que si el programa se ejecuta muchas veces el costo promedio deberá tender a la solución analítica de $51,840.




6. Un Sistema de Inventarios.
=============================

Los sistemas de inventarios contienen gran cantidad de variables que interactuan entre sí. Es usual encontrar soluciones matemáticas a los problemas de inventarios determinísticos, sin embargo cuando algunas de estas variables se tornan estocásticas el modelo matemático se vuelve considerablemente complejo o imposible de construir.

Usualmente los modelos de inventarios determinísticos asumen una demanda conocida y constante, una capacidad de producción infinita, no se permiten faltantes y los costos de ordenar y transportar el inventario tienen un comportamiento lineal.

En la realidad todas estas variables son mucho más complejas. La demanda es aleatoria y puede comportarse diferente de acuerdo a la época del año. El tiempo de entrega que transcurre entre el solicitar una orden y el momento de entrega no siempre es igual. Hay costos por mantener inventarios y costos por no tener suficiente inventario si un cliente lo solicita. Además todos estos parámetros pueden ser no lineales.

A continuación se presenta un estudio de un sistema de inventarios sencillo, para el cual es necesario identificar dos variables primordiales. La cantidad de la orden denotada por "q" y el punto de reorden denotado por "r". Por consiguiente para evaluar el funcionamiento del sistema se toman valores de "q" y "r" que minimicen los costos totales de operación.


Al inicio de cada día se revisa el nivel de inventario. Cuando el nivel de inventario está por debajo del valor de reorden de "r". Se debe solicitar un nuevo pedido. El tiempo de entrega en días del pedido está dado por:


--------------------
  Días  Probabilidad
--------------------
    1       0.30
    2       0.40
    3       0.30
--------------------


Para estimar estos tiempos en nuestro modelo utilizamos también el método de la transformada inversa.


----------------------
      r         Tiempo
----------------------
[0.00 , 0.30[     1
[0.30 , 0.70[     2
[0.70 , 1.00[     3
----------------------



La demanda se suele construir a través de datos históricos de la empresa. Estos datos se pueden usar directamente como una tabla o bien se puede tratar de construir una distribución de probabilidad que se ajuste a los datos. 

En este ejemplo se supondrá que la demanda diaria se comporta como una distribución Poisson con un parámetro de lambda de 7.



           x     _ \
         \         /\
         /\   *  e
p(x)  =  -------------    para x = 0,1,2,3,...
               x!




               
          7^x * e^(-7)
p(x)  =   ------------    para x = 0,1,2,3,...
              x!




Para la construcción de los valores se utilizará el método de la transformada inversa, de manera tal que que se construirá una tabla con los posibles valores de la distribucion poisson y su acumulada. Para construir la tabla se acumulan valores hasta superar el 0.99 de los datos.


----------------------
  x    p(x)     P(x)
----------------------
  0   0.0009   0.0009
  1   0.0064   0.0073
  2   0.0223   0.0296
  3   0.0521   0.0818
  4   0.0912   0.1730
  5   0.1277   0.3007
  6   0.1490   0.4497
  7   0.1490   0.5987
  8   0.1304   0.7291
  9   0.1014   0.8305
 10   0.0710   0.9015
 11   0.0452   0.9467
 12   0.0263   0.9730
 13   0.0142   0.9872
 14   0.0071   0.9943
----------------------



De esta forma para generar los números se utiliza el algoritmo dado por la siguiente tabla.


-----------------------------
                         -1
      r            x =  P (r)
-----------------------------
[0.0000 , 0.0009[    0
[0.0009 , 0.0073[    1
[0.0073 , 0.0296[    2
[0.0296 , 0.0818[    3
[0.0818 , 0.1730[    4
[0.1730 , 0.3007[    5
[0.3007 , 0.4497[    6
[0.4497 , 0.5987[    7
[0.5987 , 0.7291[    8
[0.7291 , 0.8305[    9
[0.8305 , 0.9015[   10
[0.9015 , 0.9467[   11
[0.9467 , 0.9730[   12
[0.9730 , 0.9872[   13
[0.9872 , 0.9943[   14
-----------------------------



La información de costos está dada por los siguientes elementos:

() El costo fijo de hacer un pedido será de $100. No importa la cantidad que se ordene se deben pagar $100 por el pedido.

() El costo del inventario se calculará con base en el inventario al final de cada día. Este cuesta $20 por unidad por día.

() Cuando algún cliente solicita un pedido y este no está disponible se estima un costo de $50 por día. Este es un costo de penalización por la insastisfacción del cliente y otros gastos.


A continuación se presenta una tabla con una ejecución de este problema de inventarios. Para la presenta tabla se han tomado los siguientes supuestos:

Inventario inicial   i = 50

Cantidad por ordenar q = 50

Punto de reorden     r = 20



---------------------------------------------
No   ini   q?   r1     t    r2       x   fin
---------------------------------------------
 1    50                   0.6825    8    42
 2    42                   0.2244    5    37
 3    37                   0.7822    9    28
 4    28                   0.4765    7    21
 5    21                   0.6172    8    13
 6    13   si  0.8245  3   0.6842    8     5
 7     5               2   0.9365   11    -6
 8    -6               1   0.5024    7   -13
 9    37                   0.2865    5    32
10    32                   0.8387   10    22
11    22                   0.6393    8    14
12    14   si  0.6911  2   0.1976    5     9
13     9               1   0.9940   14    -5
14    45                   0.8491   10    35
15    35                   0.3476    6    29
16    29                   0.1987    5    24
17    24                   0.7683    9    15
18    15   si  0.2134  1   0.2980    5    10
19    60                   0.5951    7    53
20    53                   0.0626    3    50
---------------------------------------------



El primer gráfico que se presenta se realiza con los inventarios iniciales de cada día. El gráfico muestra el comportamiento del inventario conforme avanzan los días. Observe cuando el nivel de inventario se hace negativo que se produce un faltante de productos. 


Inventario                                                 .
   |                                                          .
 50| .
   |    .                                   .
   |       .                 .
   |                            .              .
   |          .                                   .
   |             .                 .                 .
   |
   |                .                 .                 .
   |                                     .
  0|___________________._______________________________________
   | 1  2  3  4  5  6  7  8  9 10 11 12 13 14 15 16 17 18 19 20
   |                      .
   |
   |
-25|



El segundo gráfico muestra los inventarios finales de cada día. Observe que en este gráfico se muestran de forma diferente la cantidad de faltantes que en el gráfico anterior.



Inventario
   |                                                       .
 50|                                                          .
   | .
   |    .
   |                         .              .
   |       .                                   .
   |          .                 .                 .
   |
   |             .                 .                 .
   |                                  .                 .
  0|________________.__________________________________________
   | 1  2  3  4  5  6  7  8  9 10 11 12 13 14 15 16 17 18 19 20
   |                   .                 .
   |
   |                      .
-25|



Ahora es importante calcular los costos en que se incurre al utilizar este modelo de inventarios. Tal como se había señalado se deben contabilizar tres costos diferentes.



() El costo fijo por orden:

Tal como se observa en la tabla se requirió hacer un pedido en tres ocasiones distintas, por cada una de ellas se deben pagar $100.


costo_fijo = 3 * 100 = 300



() El costo por mantener inventario.

Existen muchas formas de contabilizar el costo del inventario. En nuestro caso se aplicará este costo únicamente a las piezas de inventario que no se han vendido. Observe además que se hay un faltante el costo de inventario de esa día es de 0.


costo_inventario   =  42 * 20
                   +  37 * 20
                   +  28 * 20
                   +  21 * 20
                   +  13 * 20
                   +   5 * 20
                   +   0 * 20 <-- no hay costo
                   +   0 * 20 <-- no hay costo
                   +  32 * 20
                   +  22 * 20
                   +  14 * 20
                   +   9 * 20
                   +   0 * 20 <-- no hay costo
                   +  35 * 20
                   +  29 * 20
                   +  24 * 20
                   +  15 * 20
                   +  10 * 20
                   +  53 * 20
                   +  50 * 20


costo_inventario   = 439 * 20  


costo_inventario   = $8780



() El costo de penalización.

Cuando un cliente no encuentra el producto en inventario se incurre en un costo de penalización. En nuestro caso se había estimado en $50 por unidad.

costo_penalización  =  50 * (6 + 13 + 5)  

costo_penalización  =  50 * 24

costo_penalización  =  $1200



() Por lo tanto el costo total de la política q = 50, r = 20 se puede estimar en:

costo_total  =  costo_fijo 
             +  costo_inventario 
             +  costo_penalización


costo_total  =  300  
             + 8780 
             + 1200


costo_total  = $ 10280


En esta caso sería muy interesante estudiar otras políticas de cantidad por ordenar y punto de reorden, es decir otros valores para q y r. En general se buscaría encontrar valores de q y r que minimicen el valor del costo total, dadas las condiciones de demanda y tiempo de reorden de este problema.




7. Sistemas de Colas.
=====================

Un sistema de colas o líneas de espera, es el efecto resultante cuando la demanda de un servicio supera la capacidad de proporcionar dicho servicio.

Dependiendo del sistema de colas que se desea simular puede ser que exista una sola fila y un solo servidor, una sola fila y múltiples servidores o bien múltiples filas y múltiples servidores. Finalmente los sistemas de colas se pueden unir, de forma tal que la salida de una cola se convierta en la entrada de otra.

De igual forma los servidores pueden ser cajeros, máquinas, semáforos, procesos, etc. Los clientes del servicio irán de acuerdo al tipo de servidor que se tenga asociado.


A continuación se muestra de forma gráfica la forma usual de una línea de espera con una fila y un servidor.




Fila     0     0     0__   _     0
--->    /|\   /|\   /|     |    /|\
        / \   / \   / \    |    / \
       ____________________|_________
          Clientes          Servidor



El análisis de colas matemático tradicional asume un conjunto de supuestos que pasamos a listar, cuando alguno de estos supuestos se modifica es casi imposible de resolver de forma analítica.


A continuación se enumeran estos supuestos para que cuando resolvamos el  problema en forma de simulación podemos establecer que supuestos se han modificado.


() El tiempo entre llegadas de los clientes y el tiempo de servicio siguen una distribución exponecial.

() Se asume también que la cantidad de clientes que llega a sistema por unidad de tiempo y la cantidad de clientes que son servidos por el sistema tiene una distribución poisson.

() La fuente que alimenta al sistema es de tamaño infinito.

() La forma como se brinda el servicio es de forma fifo (first in, first out), el primero en llegar es el primero en ser atendido.

() Los clientes que llegan al sistema se unen al final de cola y estos nunca se retiran una vez que han entrado en la línea de espera, es decir permancen en ella hasta ser servidos.

() Los clientes llegan en forma individual, a lo sumo llega uno por unidad de tiempo, por lo que se usan unidades de tiempo muy pequeñas.

() De igual forma solamente un cliente puede abandonar el sistema por unidad de tiempo.


() El sistema ha alcanzado un estado estable.




A continuación un ejemplo de simulación de una línea de espera típíca.


Una fábrica de electrodomésticos nacional es abastecida desde un almacén central. La mercancía llega a esta fábrica durante la noche. El personal encargado de descargar la mercancía consiste en tres personas, las cuales trabajan en grupo para descargar los camiones que llegan. Trabajan un turno de 8 horas de 12:00am a 8:00am. Estas personas pueden empezar a tomar sus alimentos a partir de las 4:00am. Si a las 4:00am se está descargando un camión, entonces se empezará a tomar los alimentos después de descargarlo.

El salario por hora que recibe este personal es de $25. El almacén recibe mercancía únicamente de 12:00am a 8:00am. Si se requiere trabajar horas extras, el salario percibido por el personal será de $37.50 por hora. Finalmente el costo de espera de un camión es de $100 por hora y el costo de tener operando el almacén es de $500 por hora.

Cuando el almacén abre sus puertas puede suceder que haya más de un camión esperando para ser descargado. Con base en la información histórica, se sabe que la distribución de probabilidad del número de camiones que están esperando al momento de que el almacén abra sus puertas es la siguiente.
 

-----------------------
Cantidad
Camiones   Probabilidad
-----------------------
   0         0.50
   1         0.25
   2         0.15
   3         0.10
-----------------------



Por otra parte, también de los registros, se ha obtenido la distribución de probabilidad del tiempo entre llegadas.

-----------------------
Tiempo
entre
llegadas
(minutos)  Probabilidad
-----------------------
   20         0.02
   25         0.09
   30         0.12
   35         0.25
   40         0.20
   45         0.15
   50         0.10
   55         0.05
   60         0.03
-----------------------



Finalmente por medio de experimentación se han obtendio las distribuciones de probabilidad del tiempo de servicio para diferentes tamaños de equipo. Tales distribuciones se muestran a continuación.


Tiempo Servicio                   Tiempo servicio
para 3 personas                   para 4 personas
-----------------------           -----------------------
Minutos    Probabilidad           Minutos    Probabilidad
-----------------------           -----------------------
  20          0.05                  20          0.05
  25          0.10                  25          0.15
  30          0.20                  30          0.20
  35          0.25                  35          0.20
  40          0.12                  40          0.15
  45          0.10                  45          0.12
  50          0.08                  50          0.08
  55          0.06                  55          0.04
  60          0.04                  60          0.01
-----------------------           -----------------------




Tiempo Servicio                   Tiempo servicio
para 5 personas                   para 6 personas
-----------------------           -----------------------
Minutos    Probabilidad           Minutos    Probabilidad
-----------------------           -----------------------
  20          0.10                  20          0.12
  25          0.18                  25          0.15
  30          0.22                  30          0.26
  35          0.18                  35          0.15
  40          0.10                  40          0.12
  45          0.08                  45          0.08
  50          0.06                  50          0.06
  55          0.05                  55          0.04
  60          0.03                  60          0.02
-----------------------           -----------------------



Con base en toda la información presentada, lo que se desea establecer es el tamaño óptimo del equipo.




Solución.
---------


En el sistema de líneas de espera que se analiza existen varios elementos estocásticos  que se deben tomar en cuenta. 

() La variable de decisión es el número de personas que se encargarán de recibir los camiones.

() El número de camiones esperando a ser servidos al iniciar el proceso.

() El proceso de llegadas de clientes al sistema y el proceso de servicio.



Iniciamos el proceso de construcción del sistema construyendo mediante el método de la transformada inversa, la distribución del número de camiones que estarán esperando al iniciar.


---------------------
      r            x
---------------------
[0.00 , 0.50[      0
[0.50 , 0.75[      1
[0.75 , 0.90[      2
[0.90 , 1.00[      3
---------------------



Utilizamos también el método de la transformada inversa para simular el tiempo entre las llegadas de los clientes.


---------------------
      r            x
---------------------
[0.00 , 0.02[     20
[0.02 , 0.10[     25
[0.10 , 0.22[     30
[0.22 , 0.47[     35
[0.47 , 0.67[     40
[0.67 , 0.82[     45
[0.82 , 0.92[     50
[0.92 , 0.97[     55
[0.97 , 1.00[     60
---------------------


Utilizamos también el método de la transformada inversa para simular el tiempo de servicio para los distintos equipos.


Para 3 personas:               Para 4 personas:
---------------------          ---------------------
      r            x                 r            x
---------------------          ---------------------
[0.00 , 0.05[     20           [0.00 , 0.05[     20
[0.05 , 0.15[     25           [0.05 , 0.20[     25
[0.15 , 0.35[     30           [0.20 , 0.40[     30
[0.35 , 0.60[     35           [0.40 , 0.60[     35
[0.60 , 0.72[     40           [0.60 , 0.75[     40
[0.72 , 0.82[     45           [0.75 , 0.87[     45
[0.82 , 0.90[     50           [0.87 , 0.95[     50
[0.90 , 0.96[     55           [0.95 , 0.99[     55
[0.96 , 1.00[     60           [0.99 , 1.00[     60
_____________________          _____________________




Para 5 personas:               Para 6 personas:
---------------------          ---------------------
      r            x                 r            x
---------------------          ---------------------
[0.00 , 0.10[     20           [0.00 , 0.12[     20
[0.10 , 0.28[     25           [0.12 , 0.27[     25
[0.28 , 0.50[     30           [0.27 , 0.53[     30
[0.50 , 0.68[     35           [0.53 , 0.68[     35
[0.68 , 0.78[     40           [0.68 , 0.80[     40
[0.78 , 0.86[     45           [0.80 , 0.88[     45
[0.86 , 0.92[     50           [0.88 , 0.94[     50
[0.92 , 0.97[     55           [0.94 , 0.98[     55
[0.97 , 1.00[     60           [0.98 , 1.00[     60
---------------------          ---------------------




La simulación de las líneas de espera se puede realizar de dos formas.

() Incrementos a tiempo fijo

() Incrementos al próximo evento.


Ambos procedimientos actualizan el estado del sistema de acuerdo a los eventos que ocurren en un intervalo específico de tiempo. La diferencia entre estos procedimientos es la manera en la cual se especifican los incrementos.


En los incrementos a tiempo fijo, como su nombre lo indica, el tiempo es adelantado siempre en intervalos que tienen una misma duración. Por ejemplo el tiempo puede ser adelantado en incrementos de 5 minutos. Cuando el reloj virtual de la simulación es adelantado en esta unidad de tiempo el sistema es actualizado al considerar todos los eventos que ocurrieron durante este período de tiempo. El sistema pasa a un nuevo estado y está listo para repetir el proceso.


En los incrementos al próximo evento, el reloj de la simulación se adelanta en intervalos que no tienen la misma duración. El tiempo de cada incremento se determina con base en el cambio más cercano que existe dentro del sistema. Por ejemplo si se sabe que van a ocurrir dos cambios importantes en el sistema uno dentro de 7 minutos y otro dentro de 11 minutos, el sistema se avanza directamente 7 minutos, al tiempo más cercano del próximo evento que está por ocurrir.



A continuación se presenta una tabla con una simulación manual del problema con un turno de operación de 3 personas. Al inicio de la tabla se ha generado un número aleatorio y existe un camión en espera de ser atendido al momento de abrir las puertas el almacén. Como la tabla es algo larga se han utilizado las siguientes abreviaciones:

Tiempo: tiempo actual del sistema
Evento: si es inicio, entrada o salida
c     : tamaño de la cola de espera
is    : inicio del servicio
ts    : tiempo del servicio
SALE  : tiempo en que sale el vehículo
te    : tiempo para la siguiente entrada
ENTRA : tiempo en que entra el vehículo
p  : tiempo de ocio del personal
l  : longitud de la cola
(i): inicio del sistema
(*): 30 minutos de almuerzo



----------------------------------------------------------------------
Tiempo  Evento  c  is      r      ts  SALE   r       te  ENTRA   p   c
----------------------------------------------------------------------
12:00   (i)     1  12:00  0.8376  50  12:50  0.4835  40  12:40   0   0
12:40   entra   2  12:50  0.1438  25   1:15  0.9897  60   1:40   0  10
12:50   sale    1   ....  ......  ..   ....  ......  ..   ....   .   .
 1:15   sale    0   ....  ......  ..   ....  ......  ..   ....   .   .
 1:40   entra   1   1:40  0.5132  35   2:15  0.0653  25   2:05  25   0
 2:05   entra   2   2:15  0.7247  45   3:00  0.4512  35   2:40   0  10
 2:15   sale    1   ....  ......  ..   ....  ......  ..   ....   .   .
 2:40   entra   2   3:00  0.0546  25   3:25  0.1548  30   3:10   0  20
 3:00   sale    1   ....  ......  ..   ....  ......  ..   ....   .   .
 3:10   entra   2   3:25  0.8460  50   4:15  0.1924  30   3:40   0  15
 3:25   sale    1   ....  ......  ..   ....  ......  ..   ....   .   .
*3:40   entra   2   4:45  0.2973  30   5:15  0.1599  30   4:10   0  65
 4:10   entra   3   5:15  0.5907  35   5:50  0.6794  45   4:55   0  65
 4:15   sale    2   ....  ......  ..   ....  ......  ..   ....   .   .
 4:55   entra   3   5:50  0.7635  45   6:35  0.9087  50   5:45   0  55
 5:15   sale    2   ....  ......  ..   ....  ......  ..   ....   .   .
 5:45   entra   3   6:35  0.2954  30   7:05  0.5899  40   6:25   0  50
 5:50   sale    2   ....  ......  ..   ....  ......  ..   ....   .   .
 6:25   entra   3   7:05  0.6195  40   7:45  0.6869  45   7:10   0  40
 6:35   sale    2   ....  ......  ..   ....  ......  ..   ....   .   .
 7:05   sale    1   ....  ......  ..   ....  ......  ..   ....   .   .
 7:10   entra   2   7:45  0.1726  30   8:15  0.7348  45   7:55   0  35
 7:45   sale    1   ....  ......  ..   ....  ......  ..   ....   .   .
 7:55   entra   2   8:15  0.1006  25   8:40  0.5542  35   8:30   0  20
 8:15   sale    1   ....  ......  ..   ....  ......  ..   ....   .   .
 8:40   sale    0   ....  ......  ..   ....  ......  ..   ....   .   .
----------------------------------------------------------------------
Total                                                           25 385
----------------------------------------------------------------------




Ahora bien, para la tabla anterior, debemos calcular el costo en que se ha incurrido.



salarios = tiempo_normal + tiempo_extra

         = 3 * 8 * $25   + 3 * (40/60) * $37.50

         = 600           + 75

         = 675



tiempo_ocioso = 3 * (25/60) * $25
              
              = 31.25



operación_almacén = (8 + 40/60) * $500

                  = 4333.3333



espera_del_camión = (385/60) * $100

                  =  641.6667

         

costo_total = salarios 
            + tiempo_ocioso 
            + operación_almacén
            + espera_del_camión


            =   675.00
            +    31.25
            +  4333.3333
            +   641.6667

            = 5681.25



A continuación se presenta una tabla con una simulación manual del problema con un turno de operación de 4 personas. Al inicio de la tabla se ha generado un número aleatorio y existe un camión en espera de ser atendido al momento de abrir las puertas el almacén. Se utilizan las mismas abreviaciones que en la tabla anterior.


----------------------------------------------------------------------
Tiempo  Evento  c  is      r      ts  SALE   r       te  ENTRA   p   c
----------------------------------------------------------------------
12:00   (i)     1  12:00  0.1140  20  12:20  0.6897  45  12:45   0   0
12:20   sale    0  .....  ......  ..   ....  ......  ..  .....   .   .
12:45   entra   1  12:45  0.6562  35   1:20  0.1847  30   1:15  25   0
 1:15   entra   2   1:20  0.9397  45   2:05  0.1470  30   1:45   0   5
 1:20   sale    1   ....  ......  ..   ....  ......  ..   ....   .   .
 1:45   entra   2   2:05  0.2256  25   2:30  0.8374  50   2:35   0  20
 2:05   sale    1   ....  ......  ..   ....  ......  ..   ....   .   .
 2:30   sale    0   ....  ......  ..   ....  ......  ..   ....   .   .
 2:35   entra   1   2:35  0.3336  25   3:00  0.1693  30   3:05   5   0
 3:00   sale    0   ....  ......  ..   ....  ......  ..   ....   .   .
 3:05   entra   1   3:05  0.0712  20   3:25  0.2036  30   3:35   5   0
 3:25   sale    0   ....  ......  ..   ....  ......  ..   ....   .   .
 3:35   entra   1   3:35  0.3748  25   4:00  0.4119  35   4:10  10   0
 4:00   sale    0   ....  ......  ..   ....  ......  ..   ....   .   .
*4:10   entra   1   4:30  0.3167  25   4:55  0.6691  40   4:50   0  20
 4:50   entra   2   4:55  0.5413  30   5:25  0.3535  35   5:25   0   5
 4:55   sale    1   ....  ......  ..   ....  ......  ..   ....   .   .
 5:25   sale    0   ....  ......  ..   ....  ......  ..   ....   .   .
 5:25   entra   1   5:25  0.6841  35   6:00  0.7998  45   6:10   0   0
 6:00   sale    0   ....  ......  ..   ....  ......  ..   ....   .   .
 6:10   entra   1   6:10  0.5232  30   6:40  0.4685  35   6:45  10   0
 6:40   sale    0   ....  ......  ..   ....  ......  ..   ....   .   .
 6:45   entra   1   6:45  0.9345  45   7:30  0.6924  45   7:35   5   0
 7:30   sale    0   ....  ......  ..   ....  ......  ..   ....   .   .
 7:35   entra   1   7:35  0.3179  25   8:00  0.0401  25   8:00   5   0
 8:00   sale    0   ....  ......  ..   ....  ......  ..   ....   .   .
 8:00   entra   1   8:00  0.8731  45   8:45  0.9379  55   8:55   0   0
----------------------------------------------------------------------
Total                                                           65  50
----------------------------------------------------------------------





Se procede a  calcular el costo de esta nueva simulación.



salarios = tiempo_normal + tiempo_extra

         = 4 * 8 * $25   + 4 * (45/60) * $37.50

         = 800           + 112.50

         = 912.50



tiempo_ocioso = 4 * (65/60) * $25
              
              = 108.3333



operación_almacén = (8 + 45/60) * $500
                 
                  = 3000



espera_del_camión = (50/60) * $100
 
                  = 83.3333



costo_total = salarios 
            + tiempo_ocioso 
            + operación_almacén
            + espera_del_camión

            =  912.50
            +  108.3333
            + 3000.00
            +   83.3333

            = 4104.16666




Si esta simulación fuera acertada, al utilizar una persona más para el servicio de descarga se ahorra una cantidad considerable de dinero. Por lo tanto resulta mejor contratar a 4 personas. En los ejercicios se realizarán simulaciones para 5 y 6 personas.




8. Resumen.
===========

En el presente capítulo se han mostrado ejemplos y aplicaciones de la simulación a diferentes ámbitos tanto científicos como industriales y de administrativos.

Iniciamos los ejemplos utilizando juegos de azar propios de los casinos, los cuales son siempre interesantes y muestran los principios básicos del proceso de simulación. Posteriormente mostramos una aplicación científica, como es el cálculo de pi, y finalmente aplicaciones a la industria y los negocios como son el manejo de los inventarios y las líneas de espera.

En general cualquier sistema que se pueda resolver de forma matemática o analítica se puede resolver mediante simulación. Sin embargo no todo sistema simulado se puede resolver de forma analítica. Sin embargo al usar esta técnica de resolución de problemas se incurre en algunos riesgos. Es importante ejecutar el proceso de simulación un número suficientemente grande de veces para que los resultados obtenidos se acerquen a la solución real del problema. Esto puede conllevar mucho tiempo humano y computacional.



 
9. Ejercicios
=============


() Ejercicio 1

Construya un programa que realice una simulación del juego de dados descrito. Asuma que se cuenta con una cantidad inicial de $20, cada vez que se tiran los dados se apuesta $1, si gana recibe una cantidad igual y si pierde debe descaontar esa cantidad de su dinero.

a. Cuál es la probabilidad de llegar a $50.

b. Cuál es la probabilidad de quedar en bancarrota.




() Ejercicio 2

El juego de casino de la ruleta consta de una rueda con 22 números. De los números 10 de ellos son de color rojo, 10 son de color negro y hay dos números de color verde, que son el cero y el doble cero. Si un jugador apuesta al color rojo y este aparece, gana la cantidad apostada. Si un jugador apuesta al color negro y este aparece, gana la cantidad apostada. Si se apuesta al color verde y este aparece, no se gana ni pierde, si se apuesta a otro color y aparece el verde se pierde.

En la mesa de juego hay dos jugadores y cada uno de ellos sigue una estrategia de juego diferente. El jugador A simplemente apuesta $1 al color rojo en cada ocasión. El jugador B inicia apostando $1 al color rojo, si gana apuesta nuevamente $1, si pierde dobla la apuesta anterior, es decir al perder por primera vez apuesta $2, si vulve a perder apuesta $4 y asi sucesivamente. Si el dinero que tiene no es suficiente apuesta todo lo que tiene.

Si cada jugador inicia el juego con $200, determine cuál de las dos estrategias es mejor.




() Ejercicio 3

Un compañía de renta de automóviles está tratando de determinar el número óptimo de autos que debe comprar. El costo promedio anual de un auto es de $75000. Además esta compañía ha recopilado las siguientes probabilidades de operación.


-----------------------           -----------------------
Autos                             Número
Rentados                          de días
por día    Probabilidad           rentados   Probabilidad
-----------------------           -----------------------
  0        0.10                      1       0.40
  1        0.10                      2       0.35
  2        0.25                      3       0.15
  3        0.30                      4       0.10
  4        0.25
-----------------------           -----------------------


Suponga que la renta diaria de un auto es de $200, el costo de no tener una auto disponible cuando es solicitado es de $150 y el costo de tener un carro sin uso es de $50. Asímismo asuma que un auto que se renta, cuando es entregado está disponible al día siguiente y que el año tiene 365 días.

a. Indique cada una de las partes del modelo de simulación involucrado.

b. Determine la cantidad de autos que óptima que debe comprar la compañía.



() Ejercicio 4

La demanda diaria de un artículo en particular está regida por una distribución binomial con parámetros de n=6 y p=0.50. El tiempo de entrega en días es una variable aleatoria poisson con un valor de lambda de 3. El costo de mantener una unidad en inventario es de $1 por día, el costo del faltante es de $10 por unidad y el costo de ordenar es de $50 por orden. Asímismo, las piezas faltantes en un ciclo son surtidas por la nueva orden cuando llega. Se desean comparar dos políticas para optimizar este proceso de inventario.

La primera consiste en ordenar cada 8 días hasta tener 30 artículos en inventario. La segunda política consiste en ordenar 30 artículos cuando el nivel de inventario es menor a 10. 

a. Indique cada una de las partes del modelo de simulación involucrado.

b. Determine mediante simulación cuál de las dos políticas es mejor.

c. Si las piezas faltantes no son sustituibles, es decir si se hace una orden y no existen piezas en inventario, esa orden se pierde, qué cambios hay que hacer al modelo y cuál política es mejor.



() Ejercicio 5

Una compañía tiene un problema de mantenimiento con cierto equipo que contiene 4 componentes electrónicos idénticos que parecen ser los causantes del problema. Este consiste en que los componentes fallan frecuentemente, forzando a que el equipo se desconecte mientras se reponen. Lo que se ha venido haciendo es reemplazar los componentes solamente cuando se descomponen. Sin embargo, existe una nueva propuesta de hacer el reemplazo de los cuatro componentes cuando falle cualquiera de ellos con el objeto de reducir la frecuencia de desconexión del equipo.

El tiempo de vida de un componente está normalmente distribuido con una media de 600 horas y una desviación estándar de 100 horas. También se sabe que es necesario desconectar el equipo 1 hora si se reemplaza un componente y 2 horas si se reemplazan los 4 componentes. Un componente nuevo cuesta $200 y se incurre en un costo de $100 por hora cada vez que se desconecta el equipo.

a. Indique cada una de las partes del modelo de simulación involucrado.

b. Determine cuál de las dos políticas anteriores es más económica. Para ello simule la operación del equipo durante unas 20000 horas.




() Ejercicio 6

Un vendedor de revistas compra mensualmente una revista el primer día de cada mes. El costo de cada ejemplar es de $1.50. La demanda de la revista en los primeros 10 días del mes tiene la siguiente distribución de probabilidad.

----------------------
Demanda   Probabilidad
----------------------
   5      0.05
   6      0.05
   7      0.10
   8      0.15
   9      0.25
  10      0.25
  11      0.15
----------------------


Al final del décimo día, el vendedor puede regresar cualquier cantidad al proveedor, quien se las pagará a $0.90 el ejmplar o bien puede comprar más revistas a $1.20 el ejemplar. La demanda en los siguientes 20 días está dada por la siguiente distribución.


----------------------
Demanda   Probabilidad
----------------------
   4      0.15
   5      0.20
   6      0.30
   7      0.20
   8      0.15
----------------------



Al final del mes, el vendedor puede regresar al proveedor las revistas que le sobren, las cuales se le pagarán a $0.60 el ejemplar. Finalmente, se asume que después de un mes ya no existe demanda por parte del público ya que para esa fecha habrá aparecido el nuevo número de la revista. Si el precio al público es de $32 por ejemplar determine la política óptima de compra.





() Ejercicio 7

La demanda diaria y el tiempo de entrega de un cierto producto siguen las siguientes distribuciones de probabilidad.


----------------------
Demanda
Diaria    Probabilidad
----------------------
   0      0.04
   1      0.06
   2      0.10
   3      0.20
   4      0.30
   5      0.18
   6      0.08
   7      0.03
   8      0.01
----------------------




----------------------
Tiempo
Entrega   Probabilidad
en Días
----------------------
   1      0.20
   2      0.30
   3      0.25
   4      0.25
----------------------



Si el producto no está disponible cuando es requerido el cliente puede esperar la llegada del nuevo lote por un tiempo limitado, es decir, si el cliente decide esperar 2 días y la mercancía no llega en ese tiempo, entonces, la demanda de ese cliente se considera perdida. La distribución de probabilidad del tiempo que un cliente está dispuesto a esperar para que se le surta su pedido es la siguiente:


----------------------
Tiempo
Espera    Probabilidad
en Días
----------------------
   0      0.40
   1      0.20
   2      0.15
   3      0.15
   4      0.10
----------------------


La información con respecto a los costos relevantes es la siguiente:
El costo de ordenar es de $100 por orden.
El costo de mantener el inventario es de $52 por unidad por año.
El costo del faltante suponiendo que el cliente decide esperar es de $20 por unidad y el costo del faltante suponiendo que el cliente decide no esperar es de $50 por unidad.

Si el inventario inicial es de 100 unidades, se debe determinar la cantidad óptima a ordenar q y el punto óptimo de reorden r. Para este ejercicio asuma que se trabajan 365 días al año.




() Ejercicio 8

Para el sistema de colas presentado, construya un programa que realice múltiples simulaciones para 3,4,5,6 personas. Indique cuál de estas opciones es mejor para la compañía.



() Ejercicio 9

Un banco emplea 3 cajeros para servir a sus clientes. Los clientes arriban de acuerdo a un proceso poisson a una media de 40 por hora. Si un cliente encuentra todos los cajeros ocupados, entonces se incorpora a una cola única que alimenta a todos los cajeros. El tiempo que dura una transacción entre un cajero y un cliente sigue una distribución uniforme entre 0 y 1 minuto. Con base en esta información determine cuál es el tiempo promedio que pasa un cliente en el sistema y cuál es la cantidad promedio de clientes en el sistema.




() Ejercicio 10

Se tiene un sistema de colas formado por dos estaciones en serie. Los clientes atendidos en la primera estación pasan en seguida a formar la cola de la segunda. En la primera estación de servicio la razón de llegadas sigue una distribución poisson con una media de 20 clientes por hora y el tiempo de servicio sigue una distribución exponencial con una media de atención de 2 minutos por persona. En la segunda estación el tiempo de servicio está uniformemente distribuido entre 1 y 2 minutos. Con base en esta información, determine cuál es el tiempo promedio que un cliente pasa en el sistema, desde que entra hasta que sale del segundo servidor. Analice además cuál de las dos colas presenta la fila de mayor tamaño.




() Ejercicio 11

Una tienda de pequeño tamaño tiene un parqueo con espacio para 6 automóviles. Los clientes llegan en forma aleatoria de acuerdo a un proceso de poisson a una razón media de 10 clientes por hora, y se van inmediatamente si no existen lugares disponibles en el estacionamiento. El tiempo que un auto permanece en el parqueo sigue una distribución uniforme entre 10 y 30 minutos.

a. Determine que porcentaje de clientes se pierde por no tener más lugares disponibles.
b. Cuál es la probabilidad de encontrar un lugar disponible en el parqueo.



() Ejercicio 12

Debido a un aumento en las ventas, cierta compañía manufacturera necesita de más espacio en su fábrica. La solución que se ha propuesto es la construcción de un nuevo depósito para almacenar los productos terminados. Este depósito estaría localizado a 30 kilómetros del lugar donde está ubicada la planta. Además de acuerdo a este nuevo plan, se requiere que al final del día se envíe al nuevo depósito la producción terminada.

Por otra parte se ha recopilado información y se sabe que la producción diaria de esta compañía tiene la siguiente distribución de probabilidad.


------------------------
Producción
diaria en
toneladas   Probabilidad
------------------------
[50 , 55[   0.10
[55 , 60[   0.15
[60 , 65[   0.30
[65 , 70[   0.35
[70 , 75[   0.00
[75 , 80[   0.08
[80 , 85[   0.02
------------------------


También se sabe que el tipo de camiones que se deben utilizar para trasladar esta producción tienen una capacidad media de carga de 5 toneladas. La cantidad de viajes que se pueden realizar cada día, en una jornada de 8 horas, depende del tiempo de carga y descarga, como también del tiempo que se requiere para recorrer los 30 kilómetros entre la planta y el depósito. Consecuentemente, la cantidad de producto terminado que un camión puede trasladar de la planta al depósito es una variable aleatoria cuya distribución de probabilidad es la siguiente:


---------------------------
Toneladas
diarias
transladadas
por camión     Probabilidad
---------------------------
[4.0 , 4.5[    0.30
[4.5 , 5.0[    0.40
[5.0 , 5.5[    0.20
[5.5 , 6.0[    0.10
---------------------------


Si la cantidad diaria producida es mayor que la cantidad que puede trasladar la flotilla de camiones, el excedente es enviado a través de otra compañía transportista a un costo de $100 por tonelada. Además el costo promedio anual de un nuevo camión es de $100000. Si se trabajan 250 días en el año, cuál es el número óptimo de camiones que esta compañía debe adquirir.



() Ejercicio 13

Una cierta compañía posee algunas máquinas en uso. El tiempo que dura en operación cada una de estas máquinas tiene la siguiente distribución de probabilidad.


-----------------------------
Tiempo entre
averías en
horas            Probabilidad
-----------------------------
[ 6.0 ,  8.0[    0.10
[ 8.0 , 10.0[    0.15
[10.0 , 12.0[    0.24
[12.0 , 14.0[    0.26
[14.0 , 16.0[    0.00
[16.0 , 18.0[    0.18
[18.0 , 20.0[    0.07
-----------------------------



El tiempo que un operador tarda en reparar una máquina tiene la siguiente distribución de probabilidad.


-----------------------------
Tiempo de
reparación
en horas         Probabilidad
-----------------------------
[ 2.0 ,  4.0[    0.15
[ 4.0 ,  6.0[    0.25
[ 6.0 ,  8.0[    0.30
[ 8.0 , 10.0[    0.20
[10.0 , 12.0[    0.10
-----------------------------


El costo de tener una máquina ociosa durante una hora es de $500 y el salario por hora para los operarios es de $50 por hora y el salario para los mecanicos que realizan la reparación es de $80 por hora. 

(a) Si se tienen 3 máquinas determine cuál es el costo de operación y mantenimiento de las mismas.

(b) Determine cuántas máquinas se deben utilizar para minimizar los costos. Pruebe con 3,4,5,6,... máquinas hasta que encuentre el número que minimiza los costos.





() Ejercicio 14

Una cadena de supermercados es abastecida por un almacén central. La mercancía que llega a este almacén es descargada en turnos nocturnos. Los camiones que se descargan en este almacén llegan en forma aleatoria de acuerdo a un proceso poisson una media de 2 camiones por hora.  El tiempo que un equipo de 3 trabajadores se tarda en descargar un camión sique una distribución uniforme entre 20 y 30 minutos. Si el número de trabajadores en el equipo se incrementa, entonces la razón de servicio se incrementa. Por ejemplo, si el equipo está formado por 4 trabajadores, el tiempo de servicio está uniformemente distribuido entre 15 y 25 minutos. Si el equipo está formado por 5 trabajadores el tiempo de servicio será uniforme entre 10 y 20 minutos. Finalmente si el equipo está formado por 6 trabajadores el tiempo de servicio tiene una distribución uniformemente distribuida entre 5 y 15 minutos. Cada trabajador recibe un salario de $25 por hora durante el turno nocturno de 8 horas. El costo de tener un camión esperando se estima en $50 por hora. Con los datos anteriores determine cuál es el tamaño óptimo del equipo de trabajo.



() Ejercicio 15

Un producto fabricado por una compañía requiere que se perforen cilindros en un bloque metálico y que se inserten pistones en ellos. 
La distribución de probabilidad de la medida del radio de los cilindros tiene la siguiente distribución:


fc(x) =  100    1.0000 <= x <= 1.0100
 


Es necesario que los pistones tengan un radio de 1.0000 pulgadas por lo menos y que se excedan lo menos posible de ese valor. La distribución de probabilidad de cuánto medirá el radio del pistón está dada por:


                -400(x - 1)
fp(x) =  400 * e              si x >= 1


El espacio libre entre el cilindro y el pistón es la diferencia de sus radios. Debido a la variabilidad del proceso ocasionalmente puede ocurrir una interferencia, es decir un espacio negativo entre un cilindro y el pistón que va adentro. 


a. Indique cada una de las partes del modelo de simulación involucrado.
b. Determine con qué probabilidad ocurre una interferencia, dadas las condiciones actuales de trabajo.


\end{verbatim}
 